\begin{Notation}\qquad
\begin{itemize}\itemsep=0pt
\item We denote by $\K$ a commutative field of characteristic zero. All the objects (vector spaces, algebras,
 coalgebras, pre-Lie algebras, $\dots$)
in this text will be taken over $\K$.
\item For any $n\in \mathbb{N}$, we denote by $[n]$ the set $\{1,\dots,n\}$.
\item For any set $\T$, we denote by $\K^\T$ the set of family $\lambda=(\lambda_t)_{t\in \T}$ of elements of $\K$
indexed by~$\T$, and we denote by $\K^{(\T)}$ the set of elements $\lambda\in \K^\T$ with a finite support.
Note that if~$\T$ is finite, then $\K^\T=\K^{(\T)}$.
\end{itemize}
\end{Notation}

\section{Typed decorated trees}

\subsection{Definition}

\begin{Definition}
Let $\D$ and $\T$ be two nonempty sets.
\begin{enumerate}\itemsep=0pt
\item A $\D$-decorated $\T$-typed forest is a triple $(F,\dec,\type)$, where
\begin{itemize}\itemsep=0pt
\item $F$ is a rooted forest. The set of its vertices is denoted by $V(F)$ and the set of its edges by $E(F)$.
\item $\dec\colon V(F)\longrightarrow \D$ is a map.
\item $\type\colon E(F)\longrightarrow \T$ is a map.
\end{itemize}
If the underlying rooted forest of $F$ is connected, we shall say that $F$ is a $\D$-decorated $\T$-typed tree.
\item If $(F,\dec_F,\type_F)$ and $(G,\dec_G,\type_G)$ are two $\D$-decorated $\T$-typed forests,
they are isomorphic if there exists a rooted forest isomorphism $f$ from $F$ to $G$ such that for any vertex $v$ of $F$,
$\dec_G(f(v))=\dec_F(v)$ and for any edge $e$ of $F$, $\type_G(f(e))=\type_F(e)$
\item For any finite set $A$, we denote by $\bfT_\T(A)$ the set of $A$-decorated
$\T$-typed trees $T$ such that $V(T)=A$
and $\dec=\id_A$, and by $\bfF_\T(A)$ the set of $A$-decorated $\T$-typed forests $F$ such that $V(F)=A$ and $\dec=\id_A$.
\item For any $n\geq 0$, we denote by $\bfT_{\D,\T}(n)$ the set of isomorphism classes of $\D$-decorated $\T$-typed trees $T$
such that $|V(T)|=n$ and by $\bfF_{\D,\T}(n)$ the set of isomorphism classes of $\D$-decorated $\T$-typed forests $F$
such that $|V(F)|=n$.
We also put
\begin{gather*}
\bfT_{\D,\T}=\bigsqcup_{n\geq 0} \bfT_{\D,\T}(n),\qquad
\bfF_{\D,\T}=\bigsqcup_{n\geq 0} \bfF_{\D,\T}(n).
\end{gather*}
\end{enumerate}
\end{Definition}

\begin{Example}
We shall represent the decorations of the vertices by letters alongside them.
If~$\T$ contains two elements, represented by $\typ{10}$ and $\typ{2}$\!, then
\begin{gather*}
\bfF_{\D,\T}(1)=\{\tdun{$d$},d\in \D\},
\\
\bfF_{\D,\T}(2)=\begin{Bmatrix}
\tdun{$a$}\tdun{$b$},\tddeux{$a$}{$b$}{10},\tddeux{$a$}{$b$}{2}, a,b\in \D
\end{Bmatrix}\!,
\\
\bfF_{\D,\T}(3)=\begin{Bmatrix}
\tdun{$a$}\tdun{$b$}\tdun{$c$},\tddeux{$a$}{$b$}{10}\tdun{$c$},\tddeux{$a$}{$b$}{2}\tdun{$c$},
\tdtroisun{$a$}{$c$}{$b$}{10}{10},\tdtroisun{$a$}{$c$}{$b$}{2}{10},\tdtroisun{$a$}{$c$}{$b$}{2}{2},
\tdtroisdeux{$a$}{$b$}{$c$}{10}{10},\tdtroisdeux{$a$}{$b$}{$c$}{2}{10},
\tdtroisdeux{$a$}{$b$}{$c$}{10}{2},\tdtroisdeux{$a$}{$b$}{$c$}{2}{2}, a,b,c\in \D\end{Bmatrix}\!.
\end{gather*}
Note that for any $a$, $b$, $c\in \D$,
\begin{gather*}
\tdun{$a$}\tdun{$b$}=\tdun{$b$}\tdun{$a$},\qquad\qquad
\tdtroisun{$a$}{$c$}{$b$}{10}{10}=\tdtroisun{$a$}{$b$}{$c$}{10}{10},\qquad\qquad
\tdtroisun{$a$}{$c$}{$b$}{10}{2}=\tdtroisun{$a$}{$b$}{$c$}{2}{10},\qquad\qquad
\tdtroisun{$a$}{$c$}{$b$}{2}{2}=\tdtroisun{$a$}{$b$}{$c$}{2}{2}.
\end{gather*}
Moreover,
\begin{gather*}
\bfF_{\D,\T}([1])=\{\tdun{$1$}\},
\\
\bfF_{\D,\T}([2])=\begin{Bmatrix}\tdun{$1$}\tdun{$2$},
\tddeux{$1$}{$2$}{10},\tddeux{$2$}{$1$}{10},\tddeux{$1$}{$2$}{2},\tddeux{$2$}{$1$}{2}\end{Bmatrix}\!,
\\
\bfF_{\D,\T}([3])=\left\{\!\!\begin{array}{l}
\tdun{$1$}\tdun{$2$}\tdun{$3$},
\tddeux{$1$}{$2$}{10}\tdun{$3$},\tddeux{$1$}{$3$}{10}\tdun{$2$},\tddeux{$2$}{$1$}{10}\tdun{$3$},
\tddeux{$2$}{$3$}{10}\tdun{$1$},\tddeux{$3$}{$1$}{10}\tdun{$2$},\tddeux{$3$}{$2$}{10}\tdun{$1$},
\\
\tddeux{$1$}{$2$}{2}\tdun{$3$},\tddeux{$1$}{$3$}{2}\tdun{$2$},\tddeux{$2$}{$1$}{2}\tdun{$3$},
\tddeux{$2$}{$3$}{2}\tdun{$1$},\tddeux{$3$}{$1$}{2}\tdun{$2$},\tddeux{$3$}{$2$}{2}\tdun{$1$},
\\
\tdtroisun{$1$}{$3$}{$2$}{10}{10},\tdtroisun{$2$}{$3$}{$1$}{10}{10},\tdtroisun{$3$}{$2$}{$1$}{10}{10},
\tdtroisun{$1$}{$3$}{$2$}{2}{10},\tdtroisun{$1$}{$2$}{$3$}{2}{10},\tdtroisun{$2$}{$3$}{$1$}{2}{10},
\tdtroisun{$2$}{$1$}{$3$}{2}{10},\tdtroisun{$3$}{$2$}{$1$}{2}{10},\tdtroisun{$3$}{$1$}{$2$}{2}{10},
\tdtroisun{$1$}{$3$}{$2$}{2}{2},\tdtroisun{$2$}{$3$}{$1$}{2}{2},
\tdtroisun{$3$}{$2$}{$1$}{2}{2},\\
\tdtroisdeux{$1$}{$2$}{$3$}{10}{10},\tdtroisdeux{$1$}{$3$}{$2$}{10}{10},\tdtroisdeux{$2$}{$1$}{$3$}{10}{10},
\tdtroisdeux{$2$}{$3$}{$1$}{10}{10},\tdtroisdeux{$3$}{$1$}{$2$}{10}{10},\tdtroisdeux{$3$}{$2$}{$1$}{10}{10},
\tdtroisdeux{$1$}{$2$}{$3$}{2}{10},\tdtroisdeux{$1$}{$3$}{$2$}{2}{10},
\tdtroisdeux{$2$}{$1$}{$3$}{2}{10},\tdtroisdeux{$2$}{$3$}{$1$}{2}{10},
\tdtroisdeux{$3$}{$1$}{$2$}{2}{10},\tdtroisdeux{$3$}{$2$}{$1$}{2}{10},
\\
\tdtroisdeux{$1$}{$2$}{$3$}{10}{2},\tdtroisdeux{$1$}{$3$}{$2$}{10}{2},
\tdtroisdeux{$2$}{$1$}{$3$}{10}{2},\tdtroisdeux{$2$}{$3$}{$1$}{10}{2},
\tdtroisdeux{$3$}{$1$}{$2$}{10}{2},\tdtroisdeux{$3$}{$2$}{$1$}{10}{2},
\tdtroisdeux{$1$}{$2$}{$3$}{2}{2},\tdtroisdeux{$1$}{$3$}{$2$}{2}{2},
\tdtroisdeux{$2$}{$1$}{$3$}{2}{2},\tdtroisdeux{$2$}{$3$}{$1$}{2}{2},
\tdtroisdeux{$3$}{$1$}{$2$}{2}{2},\tdtroisdeux{$3$}{$2$}{$1$}{2}{2}
\end{array}\!\!\right\}\!.
\end{gather*}
\end{Example}

\begin{Remark}
If $|\T|=1$, all the edges of elements of $\bfF_{\D,\T}$ have the same type:
we work with $\D$-decorated rooted forests.
In this case, we shall omit $\T$ in the indices describing the forests, trees, spaces we are considering.
\end{Remark}

